\chapter{Preparing to Debug}\label{ch:gdb-prepare}

Almost everyone debugs by inserting print statements. It is universal, it
works, and it is also why a great many people spend an afternoon on a bug that
a debugger would have located in ninety seconds. The cost is invisible
precisely because you never see the faster path.

This chapter is about getting into position: what to compile, how to start, and
how to ask GDB for help once you are there.

\section{Notation for Part IV}\label{sec:gdb-notation}

Two kinds of transcript appear in Part~\ref{part:gdb}. Shell transcripts keep
the convention of \autoref{sec:sh-notation}; GDB sessions use GDB's own prompt:

\begin{gdbcode}
(gdb) break main          # a line beginning "(gdb)" is what YOU type
Breakpoint 1 at 0x1149: file main.c, line 5.
(gdb) run                 # everything else is GDB's output
Starting program: /home/konyi/demo/app
\end{gdbcode}

C source appears with line numbers, so that breakpoints can refer to them:

\begin{ccode}
int main(void) {
    int *p = NULL;
    return *p;            /* this is line 3 */
}
\end{ccode}

\begin{keys}[0.26]
	\cmd{(gdb)}          & GDB's prompt. Never type it                              \\
	\cmd{\$1}, \cmd{\$2} & GDB's value history --- results it numbered for you       \\
	\cmd{\$rip}, \cmd{\$sp} & A machine register                                     \\
	\cmd{\$myvar}        & A convenience variable you created                        \\
	\texttt{\{n\}}, \texttt{[n]} & Placeholders, as in \autoref{sec:sh-notation}      \\
\end{keys}

\begin{note}
	The dollar sign means something different again here. In
	Part~\ref{part:shell} it introduced a shell variable, in
	Part~\ref{part:make} a make variable, and in GDB it marks a register, a
	convenience variable, or an entry in the value history. Each part's tools are
	unrelated; the character is just popular.
\end{note}

\section{Compiling with \texorpdfstring{\cmd{-g}}{-g}}\label{sec:gdb-g}

\begin{shell}
$ gcc -g -O0 -Wall -o app main.c parser.c
\end{shell}

\begin{keys}[0.20]
	\cmd{-g}    & Emit debug information: line numbers, variable names, types  \\
	\cmd{-g3}   & \dots{} and macro definitions, so \cmd{print} understands them \\
	\cmd{-O0}   & Do not optimise                                              \\
	\cmd{-Og}   & Optimise, but only in ways that keep debugging sensible      \\
\end{keys}

\begin{definition}[Debug information]
	A section of the binary (DWARF, on Linux) mapping machine addresses back to
	source: which line produced which instruction, where each variable lives, and
	what every type looks like. Without it GDB can still run the program and read
	memory, but it cannot name anything.
\end{definition}

Without \cmd{-g}, a backtrace is addresses:

\begin{gdbcode}
#0  0x0000555555555149 in ?? ()
#1  0x00007ffff7db0d90 in ?? () from /lib/x86_64-linux-gnu/libc.so.6
\end{gdbcode}

With it, the same backtrace is a story:

\begin{gdbcode}
#0  parse_line (s=0x0) at parser.c:42
#1  0x000055555555520a in main (argc=2, argv=0x7fffffffe3c8) at main.c:17
\end{gdbcode}

\begin{moral}
	\cmd{-g} costs nothing at run time. It makes the binary larger, and that is
	all --- it does not slow the program down and it does not change what the code
	does. Compile \emph{everything} with \cmd{-g}, including release builds, and
	strip the symbols later if you must ship something small.
\end{moral}

\begin{note}[Why \cmd{-O0} matters]
	Optimisation is what makes a debugger confusing. At \cmd{-O2} the compiler
	reorders statements, keeps variables in registers, merges identical code paths
	and inlines small functions out of existence. GDB then reports things that are
	true but useless:
	\begin{keys}[0.34]
		\texttt{<optimized out>}          & The variable exists in the source but nowhere at run time \\
		Stepping jumps backwards          & Instructions from several lines were interleaved           \\
		A breakpoint never fires          & That code was inlined, or folded into another branch       \\
		\texttt{\dots (inlined)} frames   & The call is real in the source and absent in the binary    \\
	\end{keys}
	\cmd{-Og} is the compromise: it applies the optimisations that do not destroy
	debug information. Use \cmd{-O0} while hunting a bug, and \cmd{-Og} for a build
	you want to be both fast and debuggable.
\end{note}

\begin{remark}
	Sometimes the bug only appears at \cmd{-O2}, and then you have no choice.
	\autoref{sec:gdb-optimised} covers reading optimised frames. But check first
	that the bug is really the optimiser's doing --- more often the optimiser has
	merely exposed undefined behaviour that was always there, and
	\autoref{sec:gdb-companions} finds it faster than GDB will.
\end{remark}

\section{Debug Symbols and Stripped Binaries}\label{sec:gdb-symbols}

\begin{keys}[0.30]
	\cmd{file app}             & Says \texttt{with debug\_info} or \texttt{stripped} \\
	\cmd{nm app | head}        & List the symbols                                    \\
	\cmd{objdump -h app | grep debug} & Which DWARF sections are present             \\
	\cmd{strip app}            & Remove the symbol table --- irreversibly            \\
	\cmd{objcopy -{}-only-keep-debug app app.debug} & Split the debug info into its own file \\
\end{keys}

\begin{shell}
$ file app
app: ELF 64-bit LSB pie executable, ..., with debug_info, not stripped
$ strip app && file app
app: ELF 64-bit LSB pie executable, ..., stripped
\end{shell}

\begin{note}
	Distribution binaries are stripped, which is why a backtrace through a system
	library shows \texttt{??}. The fix is to install the matching debug package
	--- \cmd{libc6-dbg} on Debian and Ubuntu, or the \texttt{debuginfod} service
	GDB will offer to use automatically. Do it once; a backtrace that names libc's
	frames is worth much more than one that does not.
\end{note}

\section{Core Dumps}\label{sec:gdb-core}

\begin{definition}[Core dump]
	A snapshot of a process's memory written when it crashes. Loading it into GDB
	gives you the state at the moment of death --- the stack, the variables, the
	faulting instruction --- without having to reproduce the crash.
\end{definition}

\begin{shell}
$ ulimit -c              # 0 means cores are disabled -- the usual default
0
$ ulimit -c unlimited    # for this shell only
$ ./app
Segmentation fault (core dumped)
$ gdb ./app core
\end{shell}

\begin{note}
	\cmd{ulimit -c unlimited} affects only the current shell and its children
	(\autoref{sec:sh-resources}). Where the file lands is set by
	\file{/proc/sys/kernel/core\_pattern}, and on a modern desktop it is usually
	handed to \texttt{systemd-coredump} rather than written to the current
	directory --- in which case \cmd{coredumpctl list} finds it and
	\cmd{coredumpctl debug} opens it in GDB directly.
\end{note}

\begin{moral}
	Core dumps are the answer to ``it crashed once, on the server, at three in the
	morning''. Enabling them costs nothing and turns an unreproducible report into
	an ordinary debugging session.
\end{moral}

\section{Ways to Start GDB}\label{sec:gdb-start}

\begin{keys}[0.34]
	\cmd{gdb ./app}               & Load the program; it is not running yet         \\
	\cmd{gdb -{}-args ./app -v in.txt} & \dots{} with the arguments it should get   \\
	\cmd{gdb ./app core}          & Post-mortem, from a core dump                   \\
	\cmd{gdb -p 1234}             & Attach to a process that is already running     \\
	\cmd{gdb -q}                  & Quiet: skip the licence banner                  \\
	\cmd{gdb -x script.gdb ./app} & Run a file of GDB commands first                \\
	\cmd{gdb -ex run -ex bt ./app} & Run commands given on the command line         \\
	\cmd{gdb -tui ./app}          & Start with the source window open (\autoref{sec:gdb-tui}) \\
\end{keys}

\begin{eg}[The one-line crash triage]
\begin{shell}
$ gdb -q -batch -ex run -ex bt --args ./app broken-input.txt
Program received signal SIGSEGV, Segmentation fault.
0x0000555555555185 in parse_line (s=0x0) at parser.c:42
42          while (*s != '\n') {
#0  parse_line (s=0x0) at parser.c:42
#1  0x00005555555551f9 in main (argc=2, argv=0x7fffffffe3b8) at main.c:17
\end{shell}
	\cmd{-batch} runs the commands and exits, so this fits in a shell script or a
	CI job. In four seconds you know the function, the line, and that the argument
	was \texttt{NULL}.
\end{eg}

\begin{note}[Attaching]
	\cmd{gdb -p} is how you investigate a process that is hung rather than crashed:
	attach, \cmd{bt}, see what it is waiting on, then \cmd{detach}. On most Linux
	systems \texttt{ptrace} is restricted to child processes, so attaching to an
	unrelated process needs \cmd{sudo}, or
	\cmd{sudo sysctl kernel.yama.ptrace\_scope=0} to lift the restriction for the
	session.
\end{note}

\section{The Prompt, Abbreviations and Help}\label{sec:gdb-ui}

GDB accepts any unambiguous prefix of a command, and the common ones have
one-letter forms that everybody uses.

\begin{keys}[0.20]
	\cmd{b}   & \cmd{break}      \\
	\cmd{r}   & \cmd{run}        \\
	\cmd{c}   & \cmd{continue}   \\
	\cmd{n}   & \cmd{next}       \\
	\cmd{s}   & \cmd{step}       \\
	\cmd{fin} & \cmd{finish}     \\
	\cmd{p}   & \cmd{print}      \\
	\cmd{bt}  & \cmd{backtrace}  \\
	\cmd{i}   & \cmd{info}       \\
	\cmd{l}   & \cmd{list}       \\
	\cmd{q}   & \cmd{quit}       \\
\end{keys}

\begin{keys}[0.26]
	\key{<CR>} on an empty line & Repeat the previous command. Holding it down is how you step \\
	\key{<Tab>}                 & Complete a command, a function name or a variable            \\
	\key{<C-c>}                 & Interrupt the running program and return to the prompt       \\
	\key{<C-d>}                 & Quit                                                         \\
	\key{<C-r>}                 & Search the command history --- the Readline bindings of \autoref{sec:sh-readline} all work \\
	\cmd{shell ls}              & Run a shell command without leaving GDB                      \\
	\cmd{!ls}                   & The same, abbreviated                                        \\
\end{keys}

Help is structured by topic, and is better than its reputation:

\begin{keys}[0.30]
	\cmd{help}                & The list of command classes                        \\
	\cmd{help breakpoints}    & Every command in one class                         \\
	\cmd{help break}          & One command, in detail                             \\
	\cmd{apropos watch}       & Search all help text                               \\
	\cmd{info}                & What GDB knows: \cmd{info locals}, \cmd{info breakpoints}, \dots \\
	\cmd{show}                & GDB's own settings: \cmd{show args}, \cmd{show language} \\
	\cmd{set}                 & Change them: \cmd{set pagination off}              \\
\end{keys}

\begin{note}
	\cmd{info} and \cmd{show} are easy to confuse and the distinction is worth
	learning: \cmd{info} asks about \emph{the program} --- its breakpoints, its
	stack, its variables. \cmd{show} asks about \emph{GDB} --- its own
	configuration. \cmd{info args} lists the current function's arguments;
	\cmd{show args} lists the arguments your program was started with.
\end{note}

\begin{tryit}
	Compile anything with \cmd{gcc -g -O0}, run \cmd{gdb -q ./app}, and type
	\cmd{start}. You are stopped on the first line of \cmd{main} with the whole of
	the rest of Part~\ref{part:gdb} available. Then \cmd{help running} and read the
	list --- it is one screen, and it is most of \autoref{ch:gdb-exec}.
\end{tryit}
